\chapter{The Divisor-Paradigm Closure Theorem}
\label{ch:closure}

\begin{chapterabstract}
Everything before this chapter certifies compositeness one way: a
planted congruence divisor. Could some algebraic mechanism --- norm
forms, cyclotomic identities, exponent families, polynomial
factorizations --- certify many complements composite wholesale,
escaping the $\ln r$-per-class accounting that walls in
Chapter~\ref{ch:limits}? This chapter answers with a theorem and its
falsification engine. We formalise construction schemes as
candidate families with certificate maps, prove that certificates
either \emph{plant} their divisor at design time or must
\emph{extract} it per instance (which is factoring), and show --- across
a concrete grammar $\mathcal L$ of planted mechanisms --- that
congruence covering is the unique economical one: value-set ceilings
kill univariate identities at any price, a first-moment cluster bound
caps every shared-variable identity family at $O(1)$ certified offsets,
isolation and orbit accounting close the window-dense frontier, and a
supply-side purchase argument closes the cheapest surviving family
$2.5\times$ over budget. The engine swept $13\,661$ enumerated and $41$
adversarially evolved instances: zero passes --- but the adversary's
first find was a $50\times$ accounting hole in our own referee, and
fixing it sharpened the theorem's economics. With mechanisms closed,
what governs the records is measured too: the entire game exponent is
an integrality gap (the fractional cover is perfect), and the greedy
point survives exact optimization over windows up to half the cover.
Three doors remain --- rounding mathematics, throughput, and two named
conditionals --- and we state each at its maximal form.
\end{chapterabstract}

\noindent\emph{Rigor marks, used throughout this chapter:} [P] proved at
this thesis's level of rigor; [C] conjectural; [V] empirically
validated input with the dataset stated. The theorem's honest status is
assembled from these marks in \S\ref{sec:verdict}.

\section{Five probes before the theorem}\label{sec:probes}

Five bounded probes, each with a predeclared termination gate, tested
the folk candidates for a better mechanism. The benchmark: covering
certifies $\sim$$4\,400$ offsets per nat of design entropy at $r=3$,
falling to $\sim$$1$ per nat at the endgame ($r\approx457$; measured
marginal $1.01$ nats per offset averaged over the last six moduli).

\begin{itemize}
\item[\textbf{P1}] \emph{Prime-power and composite moduli.} Terminated,
strictly dominated: the class $q\equiv a\ (\mathrm{mod}\ p^2)$ is a
subset of the class mod $p$ already killed, at double the cost
[V: $p=7$ gives $227\subset1\,613$ kills]. Composite moduli reduce to
their prime factors with only losses.
\item[\textbf{P2}] \emph{Polynomial value-set identities.} $\N=m^k-c$
makes $\N-q=m^k-d^k$ algebraically composite for every offset
$q=d^k-c$. Terminated, dominated --- but the closest call, and the
seed of Lemma~1: the best quadratic family kills $113$ offsets below
$10^5$ [V: $c=398$], yet confining $\N$ to a $k$-th-power family
surrenders $(1-1/k)\ln B$ nats, and \emph{the value set of a
degree-$k$ form below $Q$ has $O(Q^{1/k})$ members} --- at most a few
hundred killable offsets at any price.
\item[\textbf{P3}] \emph{Sierpi\'nski--Riesel exponent families.}
$\N_n=j\cdot2^n+c$: for any fixed candidate the induced constraint is
one congruence class --- ordinary covering. The exponent structure
correlates residues \emph{across} the family, a legitimate
search-throughput trick, not an entropy mechanism. Terminated.
\item[\textbf{P4}] \emph{Norm forms and splitting conditions.} ``$\N-q$
is a norm in $K$'' does not certify compositeness (primes with the
right splitting are norms too), and any usable splitting condition on
$q$ is a congruence condition mod discriminant-related moduli ---
covering in Galois clothing. Terminated.
\item[\textbf{P5}] \emph{The hidden-bias loophole.} Could constructed
$\N$ make residual complements composite more often than the boost
predicts? The banked campaigns answer: across three covers and
$1.2\cdot10^9$ scanned candidates, $|\hat\Ee-\Ee|\le0.1$ nats [V].
Any unexploited structure is worth $\le0.1$ nats. Terminated.
\end{itemize}

The probes sharpen the requirement. A paradigm-breaking mechanism must
\emph{simultaneously} (1) certify $\gg\sqrt Q$ distinct prime offsets,
(2) cost $o(\ln r)$ design entropy per offset at the margin, and (3) be
invisible in banked density data --- activating only on constructed
$\N$. No known structure does all three. The rest of the chapter turns
that requirement into a theorem about a language.

\section{Schemes, certificates, and the theorem}\label{sec:formalism}

Fix game parameters: bound $B$ (here $10^{200}$), threshold $Q$ (here
up to $2\cdot10^5$), $\ln\N\approx\ln B$.

\begin{definition}[Construction scheme]
A \emph{construction scheme} is a pair $(F,C)$:
\begin{itemize}
\item a \emph{candidate family} $F$: a parameter set
$\Theta\subseteq\mathbb Z^m$ with polynomial-size description and an
evaluation map $\N:\Theta\to2\mathbb Z\cap[1,B)$ computable in
$\mathrm{poly}(\log B)$; its \emph{design cost} is
$\mathrm{cost}(F)=\ln B-\ln|\N(\Theta)|$, the nats of search freedom
surrendered by joining the family;
\item a \emph{certificate map} $C$: on input $(\theta,q)$, $q$ prime
$<Q$, outputs $\bot$ or a divisor witness $D(\theta,q)$ with
$1<D<\N(\theta)-q$ and $D\mid\N(\theta)-q$, computable and verifiable
in $\mathrm{poly}(\log B)$.
\end{itemize}
$C$ is \emph{planted} when the divisibility holds as an algebraic
identity over the family (checkable symbolically); \emph{extracted}
when $C$ discovers the divisor per instance. The \emph{killed set} is
$K(\theta)=\{q:C(\theta,q)\ne\bot\}$, the \emph{residual set}
$U(\theta)$ its complement in the primes below $Q$, and the
\emph{effective exponent} $E_{\mathrm{eff}}(F,C)$ is the infimum, over
polynomial-time searches, of the log of expected work per hit divided
by work per candidate. For the classic cover,
$E_{\mathrm{eff}}=\Ee=|U|\cdot\mathrm{boost}/\ln\N$
[V: \S\ref{sec:validation}].
\end{definition}

The \emph{covering frontier} $\Ee^*(Q,B)$ is the minimum $\Ee$ over
congruence covers at the boost-versus-kills equilibrium of
\S\ref{sec:regimes} (e.g.\ $\Ee^*(122\,000,10^{199})=20.83$;
$\Ee^*(200\,000,\sim10^{200})\approx37.7$).

\begin{theorem}[Divisor-Paradigm Closure --- target statement {[C]},
proved on the grammar {[P]}]\label{thm:closure}
Let $(F,C)$ be a construction scheme whose certificate map is planted
and expressible in the schema language $\mathcal L$ of
\S\ref{sec:grammar}. Assume:
\begin{itemize}
\item[\emph{(H1)}] Hardy--Littlewood-type independence for shifted
primes in families: conditional on all divisibility data by primes
$<Q$ (and the sieve depth), the primality indicators of distinct
residual complements are independent with the Mertens-adjusted rates;
\item[\emph{(H2)}] no compositeness test with constant soundness on
random rough inputs runs in $o$(one modular exponentiation).
\end{itemize}
Then $E_{\mathrm{eff}}(F,C)\ge\Ee^*(Q,B)-o(1)$. Moreover, extracted
certificates on any window-dense family of targets are
factoring-equivalent, and certificate-free detection is already
employed at its known-optimal cost (one Fermat exponentiation) by the
standard engine.
\end{theorem}

Informally: \emph{within $\mathcal L$, nothing beats covering; outside
planted-$\mathcal L$, a scheme must either factor or test, and both sit
at known cost floors.} Before the lemmas, a sanity audit of everything
outside the formalism, from the million-desert study
(\S\ref{sec:million}): perfect powers cannot help (consecutive squares
near $10^{200}$ are $\sim10^{100}$ apart, so the window contains none);
hiding the desert in a genuine prime gap needs merit $\approx2\,180$
against $\approx42$ ever observed (\S\ref{sec:benchmarks}); composite
moduli are P1. The formalism is aimed at the only survivors: algebraic
identities and per-instance discovery.

\section{The schema language $\mathcal L$}\label{sec:grammar}

The falsifiable core of the theorem is a concrete grammar over
certificate mechanisms; each compiles to an offset enumerator, a
symbolic divisor witness, and a design-cost term:
\begin{align*}
\mathcal L\ ::=\ &\textsf{Congruence}(d,a) &&\N\equiv a\ (\mathrm{mod}\ d)\\
\mid\ &\textsf{PolyImage}(f\in\mathbb Z[m],\deg\le4) &&\N-q=A(m)\cdot B(m)\ \text{identically}\\
\mid\ &\textsf{ExpFamily}(\N=j\cdot a^n+c) &&\text{per-offset multiplicative-order conditions}\\
\mid\ &\textsf{MultiPoly}(A,B\in\mathbb Z[u_1..u_v],\deg\le2,v\le3) &&\N-q=A(u)\cdot B(u)\\
\mid\ &\textsf{NormForm}(K,\ \mathrm{disc}\le10^4) &&\N-q=\mathrm{Norm}_K(\alpha),\ \text{exhibited factorization}\\
\mid\ &\textsf{Cyclotomic}(\Phi_n,\ n\le200) &&\text{cyclotomic/Aurifeuillian splittings}\\
\mid\ &\textsf{Compose}(\mathcal L,\mathcal L) &&\text{depth}\le2.
\end{align*}
Excluded by construction: ``certificates'' that need a primality
decision to define membership (circular), and per-instance divisor
discovery --- the extraction horn, handled by the theorem's second
clause rather than by the grammar. Grammar extensions are a human act
with a hardening rule (\S\ref{sec:engine7}): any extension to $v\ge4$
shared variables must ship with the dense-frontier analysis of
Lemma~1(c) re-proved first.

\section{Lemma 1: entropy accounting for planted certificates}\label{sec:lemma1}

\emph{Statement.} For planted $(F,C)$, the amortised design cost per
certified offset is at least $\ln(1/\delta_{\mathrm{eff}})-O(\ln\ln B)$,
where $\delta_{\mathrm{eff}}$ is the density of integers in a generic
window of width $Q$ at height $B$ for which the scheme can exhibit its
witness with polynomial-range parameters; and for every family in
$\mathcal L$ this is no cheaper than covering. The proof is a case
sweep, and the cases are the chapter's real content.

\paragraph{(a) Congruence classes [P].} $\delta_{\mathrm{eff}}=1/d$,
cost $\ln d$: the covering accounting itself, with composite $d$
dominated by its prime factors (P1).

\paragraph{(b) Univariate identities [P].} A degree-$k$ identity's
certified offsets lie in the value set of a degree-$k$ form below $Q$:
at most $O(Q^{1/k})$ of them, \emph{at any price}. At $Q=122\,000$:
$349$ ($k{=}2$), $50$ ($k{=}3$), $19$ ($k{=}4$), $7$ ($k{=}6$); and
$\Phi_n$ has degree $\varphi(n)\ge2$ for $n\ge3$, so
\textsf{Cyclotomic} sits inside the same ceiling. Both families are
closed \emph{price-free}: the engine's pass bar (below) demands
$>3\sqrt Q\approx1\,048$ distinct offsets, and $349<1\,048$.

\paragraph{(c) Multivariate identities --- the hard case.} Here the
representable set can be window-dense (products of two binary quadratic
forms reach $\sim$$9\,300$ of the $17\,983$ prime offsets at
$Q=200\,000$ [V]), so existence arguments must be replaced by
\emph{access} arguments. Write the factors' degrees $d_1,d_2\ge2$
(degree-1 factors are sieve-equivalent), $d=d_1+d_2$, $v$ shared
variables, and $\alpha=\min(1,v/d)$ the value-set density exponent
[V: measured $\hat\alpha=0.761$ vs $0.750$ predicted for $(2,2,3)$].

\emph{Cluster ceiling, $v<d$} [P modulo standard equidistribution of
form values]. Distinct values of $G=A\cdot B$ below $X$ number
$\sim X^\alpha$; a first-moment count of $k$-clusters inside width-$Q$
windows gives $X^\alpha(X^{\alpha-1}Q)^{k-1}$, so clusters exist only
up to
\[
k_{\max}(\alpha)=1+\Bigl\lfloor\frac{\alpha\ln X}{(1-\alpha)\ln X-\ln Q}\Bigr\rfloor,
\]
finite iff $\alpha<1-\ln Q/\ln X$ (threshold $\alpha^*=0.9735$ at
$X=10^{200}$, $Q=2\cdot10^5$). Every $v<d$ combination has
$\alpha\le(d-1)/d\le0.9<\alpha^*$:

\begin{table}[h]
\centering\small
\caption{The cluster ceiling: maximum certifiable offsets per candidate
for shared-variable identity families $(d_1,d_2,v)$ at $X=10^{200}$,
$Q=2\cdot10^5$. A scheme needs $\sim$$18\,000$.}\label{tab:kmax}
\begin{tabular}{lrr@{\qquad}lrr}
\toprule
combo & $\alpha$ & $k_{\max}$ & combo & $\alpha$ & $k_{\max}$\\
\midrule
$(2,2,2)$ & 0.500 & 2 & $(2,2,3)$ & 0.750 & 4\\
$(2,3,3)$ & 0.600 & 2 & $(2,3,4)$ & 0.800 & 5\\
$(2,4,4),(3,3,4)$ & 0.667 & 3 & $(2,4,5),(3,3,5)$ & 0.833 & 6\\
$(3,4,5)$ & 0.714 & 3 & any $v<d$ & $\le0.9$ & $\le13$\\
\bottomrule
\end{tabular}
\end{table}

For the flagship case $(2,2,2)$ this is existence, not hardness: at
$X=10^{200}$ the expected number of 3-clusters is
$Q^2X^{-1/2}\approx10^{-89}$, so a same-variable quartic-product scheme
certifies \emph{at most two} offsets per candidate with high
probability over any construction ensemble --- versus $18\,000$ needed.
$\delta_{\mathrm{eff}}=\Theta(X^{-1/2})$ is an existence statement,
with the measured enumeration wall (work per certified in-window offset
scaling as $\mathrm{height}^{\theta}$, $\theta$ measured
$0.50$--$0.74$ against the theoretical $1/2$ [V]; windows at height
$10^{10}$ already contain zero enumerable hits at width $10^4$) as its
algorithmic shadow --- $\ge10^{100}$ operations per offset extrapolated
to $B=10^{200}$.

\emph{The dense frontier, $v\ge d$} [P at draft rigor]. When the value
set is window-dense, two elementary propositions close access anyway.
\emph{Isolation}: around a solution $u_0$ with $G(u_0)$ in the window,
any lattice displacement $\delta$ keeping $G(u_0+\delta)$ within a
width-$2Q$ window must satisfy $|\nabla G\cdot\delta|\lesssim Q$ and
$|\delta^{\mathsf T}H\delta|\lesssim Q$; the solution slab has
principal extents $Q\,X^{-(d-1)/d}$ and $Q^{1/2}X^{-(d-2)/(2d)}$ ---
at $d=4$, $X=10^{200}$: $10^{-145}$ and $10^{-47}$. Every extent is
below $\tfrac12$, so the only lattice point is $\delta=0$:
\emph{window solutions are pairwise isolated; there is no local
navigation between certified offsets --- dense existence is
non-constructive by geometry, not by hardness.} \emph{Orbit
accounting}: an infinite group $\Gamma$ of structured moves either
preserves $G$ (one offset per orbit), or preserves a factor --- so every
reachable value carries the planted divisor $a_0=A(u_0)$: enumerable
$a_0$ is a congruence certificate (covering economics), balanced $a_0$
is back in the isolation sliver --- or mixes generators of both
automorphism groups, whose coarse log-lattice of reachable factor pairs
holds $\sim(\log X)^2\approx2\cdot10^5$ points against a target of
relative width $10^{-195}$: expected hits $10^{-190}$. Any remaining
access route must solve $G(u)=m$ for \emph{prescribed} $m$ --- but a
product-form representation of $m$ \emph{is} a nontrivial factorization
of $m$, so a prescribed-value solver factors the scheme's own rough
complements: the extraction horn, not a new hypothesis. Disjoint-variable
products reduce directly to covering (the free small factor is the
certificate; \textsf{NormForm} reduces to \textsf{MultiPoly} through
exactly this route).

\paragraph{ExpFamily: closed by supply, not price [P, quantitative].}
Order conditions are the cheapest certificates in the grammar: $s$
divides $j\cdot a^n+c-q$ along an arithmetic progression of $n$ when
the discrete logarithm of $(q-c)j^{-1}$ lands in
$\langle a\rangle\bmod s$, at design cost $\ln\mathrm{ord}_s(a)$ ---
as low as $\ln3=1.10$ nats. But $s\mid a^{\mathrm{ord}}-1$ forces
$a^{\mathrm{ord}}>s$, so cheap certificates need small orders in large
$s$, and such $s$ supply only $\sim\pi(Q)\,\mathrm{ord}/s$ eligible
offsets each [e.g.\ $s=2801$, order 5: about 20]. Solving the resulting
purchase problem exactly --- sort all $s<10^5$ by price, buy cheapest
first until the pass bar of $1\,048$ offsets --- costs $1\,151$ nats
for every base $a\in\{2,3,5,7\}$, against the entire
$\ln B=460$-nat budget; the best \emph{budget-feasible} design
certifies $419$ offsets. The strongest non-covering family in the
grammar is $2.5\times$ over budget at the bar, ceiling $419$ against
$18\,000$.

\paragraph{Compose: empty at scale [P].} Two nontrivial family
constraints on the same $\N$ are contradictory: $\N=m^2-c_1=m'^2-c_2$
forces $(m-m')(m+m')=c_1-c_2$, so $m\le|c_1-c_2|$, while
$\N\sim10^{200}$ needs $m\sim10^{100}$ [V: ordinary shift pairs have
$O(1)$ solutions, largest $m=1073$ in the swept instances]. The only
surviving composition is family-plus-congruence --- a capped family
riding on ordinary covering, still under the bar. This explains
structurally the additivity the engine measured (no super-additive
interaction in any composed sample).

\paragraph{The economics, restated.} The generative tier
(\S\ref{sec:engine7}) forced a sharpening: with every mechanism charged
for \emph{all} offsets it claims and capped by the budget, a bulk
mechanism must deliver offsets at
$<\ln B/(3\sqrt Q)\approx0.44$ nats each --- and the honest comparison
is not covering's \emph{marginal} endgame ($1.01$ nats/offset) but its
\emph{average}: $437$ nats buy $10\,607$ certified offsets, $0.041$
nats each. Covering wins not at the margin but on average, because
$\ln r$ nats buy $\sim\pi(Q)/(r{-}1)$ offsets at once, whereas every
known algebraic mechanism buys $O(1)$ offsets per unit of entropy.

\section{Lemma 2: correlation exhaustion}\label{sec:lemma2}

\begin{proposition}[Shared divisors are small {[P]}]\label{prop:shared}
Let $\N$ be even, $q_1\ne q_2$ odd primes $<Q$, and $d>1$ with
$d\mid\N-q_1$ and $d\mid\N-q_2$. Then $d\mid q_1-q_2$, so every prime
factor of $d$ is $<Q$.
\end{proposition}
\begin{proof}
Subtract.
\end{proof}

Consequence: every divisor event coupling two complements is measurable
in the $\sigma$-algebra generated by small-prime divisibility --- which
is \emph{exactly} the information the cover fixes (classes) and the
sieve harvests ($k$-residues), and the engine consumes all of it: any
``$k$-side'' congruence trick selects a sub-progression the sieve
already identifies, and the survivor count is a sufficient statistic
(\S\ref{sec:engine}). Under H1, the residual lottery then factorizes,
$E_{\mathrm{eff}}$ decomposes as design cost plus $\sum_qp_q$, and that
is the covering objective [P conditional on H1]. Empirical support for
H1 at our scales is P5's $|\hat\Ee-\Ee|\le0.1$ over $1.2\cdot10^9$
candidates [V]; the corollary worth stating is that
\emph{luck-shaping} --- choosing $\N$ to make residual events
favourably dependent --- has no purchase beyond $0.1$ nats, because
dependence requires shared small divisors, which is covering by
definition.

\section{Lemma 3: the detector floor, posed as an open problem}\label{sec:lemma3}

The last escape is orthogonal to certificates: make \emph{testing}
cheaper than a modular exponentiation.

\begin{openproblem}[H2: sub-modexp compositeness detection on rough inputs]
\label{op:H2}
Let $n$ be uniformly random among integers in $[X,2X]$ with no prime
factor below $(\log X)^{O(1)}$, and let $M(b)$ be the bit cost of
$b$-bit multiplication. Does a randomized algorithm exist that decides
compositeness of $n$ with soundness $\ge2/3$ (one-sided; ``don't
know'' allowed) in time $o(M(\log X)\log X)$ --- asymptotically cheaper
than one modular exponentiation? Variants: (i) constant recall $r$ at
cost $c\cdot$modexp with $c/r<1$; (ii) amortized: certify a constant
fraction of the composites among $n_1,\dots,n_k$ in $o(k)$ modexps
total, beyond the small-divisor mass reachable by product
trees~\cite{Bernstein}; (iii) with polynomial advice depending on $X$
only.
\end{openproblem}

Known art sits firmly at the floor: trial division and batch gcd have
recall zero on rough inputs by construction; Jacobi symbols
equidistribute (soundness $o(1)$); Fermat, Miller--Rabin, and Frobenius
tests all cost $\Theta(M\log X)$; and no lower bound is known in any
realistic model --- even a conditional one in a restricted arithmetic
model would be new. The problem has a live industrial consumer: the
record engine spends $\ge95\%$ of its time on exactly this primitive
($\sim$17 tests per candidate, $e^{20.8}$--$e^{37.7}$ candidates per
campaign), as do the large distributed prime hunts; a detector at cost
$c$ and recall $r$ multiplies throughput by $1/(c+(1-r))$, and record
difficulty is $\log(\text{compute})$. A positive answer breaks the cost
floor of a whole genre of searches --- which is the more interesting
outcome, and why the problem is posed rather than assumed.

\section{The falsification engine}\label{sec:engine7}

A theorem whose proof is a case sweep earns its keep by inviting
attack. The engine's contract: a schema instance \emph{passes} ---
i.e.\ threatens the theorem --- iff simultaneously (1) it certifies
$\ge3\sqrt Q$ distinct prime offsets below $Q$ (clears the univariate
ceiling), with factor nontriviality spot-checked; (2) its cost per
offset, charged for \emph{every} offset claimed and budget-capped,
beats the covering benchmark; (3) its predicted effect would be
invisible in the banked density data (P5-consistency). Any pass is
audited by hand before it counts. The scorer re-derives every number
from primitives --- deterministic sieve, the pinned residual set,
symbolic factorization --- and \emph{proposals are data, never code};
fitness is capped, monotone in verified counts only, with no reward for
prose, novelty, or runtime.

\paragraph{Tier 0 (calibration).} The engine reproduces every P1--P4
hand number exactly, and surfaced two refinements worth keeping: the
covering endgame is $\approx1.0$ nats/offset at $Q=122\,000$ (an
earlier $2.5$-nat figure was context-bound), and the binding constraint
crosses over from budget to equilibrium between $Q=122\,000$ and
$Q=200\,000$ --- both now part of \S\ref{sec:regimes}.

\paragraph{The exhaustive tier.} $13\,661$ schema instances across the
full grammar: \textbf{zero passes}. Verdicts: $400$
congruence-composite instances scored \textsc{baseline} (inside the
paradigm --- the paradigm cannot escape itself); $12\,579$
\textsf{PolyImage} rejected on the value-set ceiling (best instance
remains $k{=}2$, $c{=}398$, $113$ offsets); $36$ \textsf{ExpFamily},
$324+4$ \textsf{MultiPoly}, $60$ \textsf{NormForm}, $198$
\textsf{Cyclotomic}, $60$ \textsf{Compose} rejected per the lemmas,
with all composition overlaps within $3\sigma$ Poisson of additive.
One methodological note kept for honesty: the first run produced $16$
spurious passes --- congruence-containing schemas leaking through the
escape filter with fresh-universe kill counts. The fix was semantic
(baseline family detection), not numeric; zero real passes before and
after. Near-miss board: \textsf{ExpFamily} at $3.1$--$4.2$ nats per
residual kill --- three times the covering endgame, exactly as P3
predicted --- then nothing below $10$.

\paragraph{The generative tier, and the referee it broke.} The
adversarial arm ran on a hosted evolutionary search service
(AlphaEvolve), under a threat model written in advance: evolved
candidates are \emph{programs}, so their only contract is to emit
schema JSON; they run sandboxed, never compute kills, never report
fitness; grading happens in a fresh interpreter. Within about ten
candidates the search drove apparent fitness $0.27\to0.46$ ---
not by finding mathematics, but by exploiting a hole in our scorer:
per-offset costs were billed only against the $867$ residuals of the
benchmark cover, so an \textsf{ExpFamily} design could advertise
$10\,416$ offsets costing $22\,900$ nats against a $460$-nat budget
and be credited anyway --- a $50\times$ overdraft. The audit protocol
caught it before any claim was made (a pass requires hand audit by
construction); the fix --- charge every claimed offset, cap coverage at
$\ln B/\mathrm{cost}$ --- re-scored the champion to $0.046$ and
\emph{strengthened} the theorem, yielding the budget constraint, the
$0.041$-nat average bar, and the purchase argument that closes
\textsf{ExpFamily} quantitatively (\S\ref{sec:lemma1}). The run then
converged: $41$ candidates, four families, zero passes, best corrected
fitness $0.0919$ --- the same near-miss family the exhaustive tier had
already ranked first, tuned to the edge of what order conditions allow,
stalled $53\times$ short of the bar.

\begin{table}[h]
\centering\small
\caption{Generative-tier empirical ceilings per family under corrected,
budget-aware accounting. ``eff.\ offsets'' is the budget-capped
certified count; the pass bar is $>1\,048$ offsets at $<0.44$
nats/offset, against covering's average $0.041$.}\label{tab:evolved}
\begin{tabular}{lrrrr}
\toprule
family & best fitness & eff.\ offsets & nats/offset & vs the bar\\
\midrule
\textsf{ExpFamily} & 0.092 & 209 & 1.10--2.20 & $53\times$ too expensive\\
\textsf{Compose} (2-deep) & 0.076 & 695 & 8.86 & $215\times$\\
\textsf{Cyclotomic} & 0.039 & 348 & 8.71 & $211\times$\\
\textsf{PolyImage} & 0.001 & 56 & 46.1 & $1118\times$\\
\bottomrule
\end{tabular}
\end{table}

Neither structural fact in Table~\ref{tab:evolved} is a search
artifact: no family clears both filter halves at once (coverage is
bought only at $8$--$9$ nats per offset; cheapness only at $O(1)$
offsets per condition), and the ordering matches Lemma~1's case
analysis exactly. The methodological finding generalises beyond this
thesis: \emph{adversarial search against a formal claim is most
valuable as a referee-hardening device, and should be reported that
way, not as a failed refutation} --- the designed behaviour of an audit
protocol in which candidates emit data, graders run fresh, and no pass
counts without human eyes.

\section{Door 1: the rounding gap}\label{sec:door1}

With mechanisms closed, the covering objective itself is the last
mathematical lever: how far from optimal are the covers? The answer
reframes the whole game.

\paragraph{The gap is everything.} The set-cover LP relaxation of the
benchmark instance ($Q=122\,000$; moduli fixed at the equilibrium set,
the $88$ odd primes $3..457$; one residue class to choose per modulus)
reports \emph{fractional coverage $100.00\%$} --- the classes carry
$\sum_r\pi(Q)/(r{-}1)\approx1.86\,\pi(Q)$ of kill mass, and fractional
routing covers every prime, at both $Q=122\,000$ and $Q=200\,000$. The
LP certifies nothing about integral assignments; what it proves is a
framing: \textbf{the entire game exponent is an integrality/rounding
gap}. Naive independent rounding leaves $e^{-1.86}\approx15.6\%$ of
primes uncovered; adaptive greedy leaves $7.56\%$; the fractional
optimum leaves $0$. Every absolute percent of residual removed is
$\approx2.7$ nats $\approx15\times$ less compute.

\paragraph{The landscape, measured.} Over the full $e^{436}$-point
assignment ensemble, random assignments have $|U|\sim(\mu=1\,377,
\sigma=21.0)$; the measured $\sigma$ sits \emph{below} the
independence prediction ($34.8$) --- negative correlations thin the
tail. A first-moment (annealed) threshold lands at $u^*\approx758$;
greedy-plus-descent achieves $867$, already $z=24.3$ of the $29.5$
available sigmas into the tail. Diagnostics: $19$ of $88$ moduli are
choice-tight; random-start coordinate ascent lands $90$--$110$ primes
\emph{worse} than greedy (adaptive construction is worth
$\approx2.4$ nats over naive local search); and pairwise class
agreement between independent solutions is $2.4\%$ --- chance level, a
fully degenerate landscape with no backbone for crossover methods to
graft.

\paragraph{Exactness, thrown at it.} The greedy point survived
everything: $41$ exact CP-SAT~\cite{ORtools} windows of $7$--$15$
moduli (warm-started, incumbent-floored) --- never beaten, always
matched, with $21$ of $22$ computed LP window bounds \emph{certifying}
window-optimality (the one uncertified window contains the smallest
moduli, the global LP's fractional looseness in miniature); exact
re-solves of the $30$ and then $44$ \emph{largest} moduli
simultaneously --- half the cover in one subproblem --- unchanged at
$867$; and belief-propagation-guided decimation (validated on a toy to
the exact optimum) landing at $874$, which its own local search cannot
improve --- a third mutually-isolated locally-exact basin
($867/874/968$), while the naive Bethe free energy slides toward the
fractional fantasy at low temperature, sharing the LP's blindness to
integrality (a survey-propagation treatment~\cite{MPZ} is the remaining
theory item). Running verdict: $867$ is at or very near the
\emph{quenched} optimum for this moduli set, the annealed $758$ likely
reflecting an annealed-versus-quenched (clustering) gap rather than
reachable headroom; the honest Door-1 prize is $0$--$2.6$ nats.

\paragraph{The open programme.} What has \emph{not} been tried is a
different ensemble: the covering constructions of Erd\H{o}s--Rankin
and FGKMT \cite{Rankin,FGKT,Maynard,FGKMT} are precisely a technology for
rounding fractional covers of primes with quantified loss --- layered
smooth-kill middle bands, second-moment control, a hypergraph-matching
endgame (cf.\ Moser--Tardos-style resampling~\cite{MT} for the
algorithmic version). Their loss is $o(1)$ in an asymptotic regime that
is not ours (their moduli reach $y/2$; ours cap at $457$), which is
exactly why this is a thesis-grade open question rather than a lookup:
\emph{determine the achievable residual fraction at
$\pi(Q)\approx10^4$ targets under budget $\theta(457)$.} Closing even
half the greedy-to-fractional gap is $e^{10}\approx22\,000\times$
compute --- the difference between a volunteer network and a
GPU-month for the $200\,000$ rung.

\section{Three doors, and the verdict}\label{sec:verdict}

Every attack generated in this programme reduces to one of three
irreducible doors; the closure argument is that there is no fourth.
\emph{(i)} Divisor information is fully harvested: $k$-side congruence
tricks select sub-progressions the sieve already identifies.
\emph{(ii)} Correlation requires a shared divisor, hence is covering
(Proposition~\ref{prop:shared}), and P5 bounds everything residual at
$0.1$ nats. \emph{(iii)} A compositeness certificate is planted (then
Lemma~1's accounting applies) or extracted (then it factors); and
certificate-free detection already runs at its conjectured floor (H2).

\begin{description}
\item[Door 1 --- rounding mathematics.] Prize $0$--$2.6$ nats on
current evidence, with the FGKMT-transfer ensemble as the genuine
unknown; every nat multiplies every future rung. Cheap, and the
covering monopoly theorem is what makes investing in it safe.
\item[Door 2 --- throughput.] The search is progress-free, stateless
per ticket, with millisecond verification: the best-shaped volunteer
workload there is, and the model on which the neighbouring
record-sports (GIMPS, PrimeGrid --- themselves covering problems in the
Sierpi\'nski--Riesel line) built everything. At the historical peak of
volunteer computing ($\sim$$2\times10^5$ of our fleets), the
$Q=200\,000$ rung falls in about two days; at $1\%$ of it, months.
The capital-intensive twin is a Montgomery-modexp ASIC rack
($\sim$$e^6$ over the fleet). The quantum door is closed at any
projected constants: Grover's halved exponent meets a
$\sim$$10^{11}$-fold oracle slowdown, crossover $\Ee^*\approx41$--$55$,
at or above every target of interest.
\item[Door 3 --- the paradigm.] This chapter. Status:
\textbf{proved on the swept grammar} [P at draft rigor] --- every
family closed by a proof of its own kind (value-set ceilings, cluster
bound, isolation plus orbit accounting, purchase argument, empty
intersection) --- \textbf{conjectural at the boundary}: H1 (supported
at $0.1$-nat resolution by $1.2\cdot10^9$ candidates), H2 (posed,
Open Problem~\ref{op:H2}), the $o(1)$ terms and equidistribution
inputs, and mechanisms outside $\mathcal L$ altogether --- neither a
planted divisor identity nor per-instance discovery --- of which none
is known.
\end{description}

The programme consequence is clean: with the covering monopoly a
result, record progress is governed by Door 1 and Door 2 --- and the
honest forecast for each game is Table~\ref{tab:walls}, unchanged by
anything a mechanism was hoped to deliver.
